\documentclass{article}
\usepackage{amsmath,hyperref}
\begin{document}
\section*{A route and information cutoff on the complete Cave01 tile graph}
\paragraph{Provenance and kernel.}
Use the exact published Cave Circuit 01 visibility graph and SDF.
This experiment treats adjacency as an \emph{authored travel
abstraction}. Each traversal or wait takes one model tick.
Graph weights from the visibility model are not travel times or
distances. The physically authorized passages, collision envelope,
time scaling, and observation law have not been established on
this distal loop. Even two near-looking placed tiles may be on
separate branches at different elevations. No field or full
three-dimensional lower bound follows.

The complete connected graph has 158 nodes, 163 edges,
19 degree-three junctions and 9 leaves. There are 101 graph
bridges and 6 independent cycles. The chosen exploration slice
is not the earlier 13-tile proved-motion subset. A distal chord
between 103 and 104 has a 19-edge alternate in the graph.
Starting at tile 103, require an inspection at tile 132 with one
service tick and return to tile 103. The two authored worlds
disable respectively the edge between 115 and 118 and the edge
between 117 and 119. Both have a ten-edge outbound and return
route when the world is known. Removing both edges separates
start from inspection on the full graph, not just a corridor
list. Only one of the two closures actually occurs.

\paragraph{Comparison: no frozen waypoint macros.}
Permit every graph walk, turn, wait, reversal and detour of the
same legal abstract unit controls. Before identification both worlds
must admit the same chosen traversal. A declared local inspection
within $r$ graph edges of any of the four endpoints discloses the
world perfectly. A delivered remote status may disclose it as
well at cut $\tau$; without it the timer is infinite. This is an
observation-model premise, not a deduction from tile visibility or
a ready-to-use transport or sensor across those distances.

The exact optimum can be computed by splitting a policy at its first
observation. Its prefix is a path using common edges. After
identification the remaining optimum is ordinary shortest path to
inspection, the service cost, and shortest path home. For local
reveal, take the shortest common prefixes to the stopping set.
For timed reveal, any vertex reachable by a non-revealing prefix
can serve at cut $\tau$, padding an early arrival with waits.
Minimize the maximum of the two matched residual costs plus the
prefix cost. Conversely every finite successful policy has such a
first split, as the joint graph cannot reach the goal. This gives
both a positive policy and a lower bound on all policies in this
finite abstraction. A Boolean horizon game independently checks
feasibility at the optimum and its failure one tick earlier.

For this graph and these controls the result is
\[
  21 + \min\bigl(\tau,\,2\max(0,4-r)\bigr).
\]
Without a report, endpoint-only observation gives 29 and adding
one sensing edge gives 27. A deadline of 24 under endpoint-only
sensing requires the report by cut 3, and cut 3 suffices. The same
deadline requires no remote report if locally reliable sensing
extends three edges. The sweep covers 240 timing, local scope and
service duration combinations. An ordinary full-graph contingent
route planner has the same optimum.

\paragraph{What would justify promoting it.}
A motion claim would require interior containment and mandatory
actual portal cuts, known body and dynamics bounds, and a route
certificate from the placed collision geometry that excludes
unmodeled chords. Contact and sensing are independent layers:
authored physical edge removal does not rewrite the native RF
graph. A separate source and finite latency must connect the
correct gate state to the chooser. The consequence of a later
fault on the same unsettled inspection and ongoing independent
work also remains untested. This is held-in graph development,
not an independent cave selection, a comparison against continuous
receding-horizon control, or I2 acceptance.
\end{document}
